\chapter{KI-gestützte Monitoring-Systeme in digitalen Informationsumgebungen}

Die theoretischen und technologischen Grundlagen digitaler Informationssysteme finden in modernen digitalen Umgebungen ihre praktische Anwendung, wobei \ac{KI}-gestützte Monitoring-Systeme als zentrale operative Einheiten fungieren. Diese Systeme operieren innerhalb von Plattformen, die nicht mehr nur als neutrale Vermittler, sondern als extraktive Apparate verstanden werden können, die darauf ausgelegt sind, menschliche Erfahrung systematisch in Verhaltensdaten zu übersetzen.\vglfootcite[21]{Srnicek2017}\vglfootcite[4]{Zuboff2019} In diesem Kontext werden die in Kapitel 2 erläuterten soziotechnischen Strukturen genutzt, um massive Volumina an Verhaltensdaten zu sammeln, zu verarbeiten und zu analysieren, was die Grundlage für Personalisierung, Empfehlungen und Vorhersageprozesse bildet.\vglfootcite[1]{Laudon2016}\vglfootcite[2]{Lyon2012} Der folgende Abschnitt verdeutlicht, wie diese Systeme die Grenze zwischen digitaler Infrastruktur und sozialer Interaktion auflösen, um Verhalten messbar und schließlich steuerbar zu machen.

\section{Datenerfassung und digitale Verhaltensanalyse}

Die digitale Verhaltensanalyse beginnt mit einer umfassenden Architektur der Datenerfassung, die jeden Aspekt der Nutzerinteraktion in ein quantifizierbares Format transformiert. Dieser Prozess der Datafizierung bildet die notwendige Voraussetzung dafür, dass Monitoring-Systeme ein kontinuierliches Abbild der Realität erstellen und darauf basierend algorithmische Entscheidungen treffen können.

\subsection{Social-Media-Plattformen}

Social-Media-Plattformen fungieren als primäre digitale Informationssysteme, in denen menschliche Kommunikation und soziale Beziehungen in maschinenlesbare Datenströme übersetzt werden.\vglfootcite[1]{VanDijck2014} Jede Interaktion -- sei es ein \enquote{Like}, ein Kommentar, das Teilen eines Beitrags oder die bloße Verweildauer auf einem Inhalt -- dient als Datenquelle für \ac{KI}-gestützte Monitoring-Systeme.\vglfootcite[12]{Narayanan2023} In Plattformen wie Facebook, Instagram oder TikTok wird das Nutzerverhalten nicht nur punktuell erfasst, sondern einer kontinuierlichen Überwachung unterzogen, die darauf abzielt, einen sogenannten Verhaltensüberschuss (Behavioral Surplus) zu generieren.\vglfootcite[4]{Zuboff2019}

Die Rolle der Künstlichen Intelligenz besteht hierbei darin, diese massiven, oft unstrukturierten Datenmengen in Echtzeit zu analysieren, um Muster in den Präferenzen und emotionalen Zuständen der Nutzer zu erkennen.\vglfootcite[1]{MetzlerGarcia2023} Beispielsweise nutzt TikTok einen rein algorithmisch gesteuerten Feed, der auf einer extrem hohen Frequenz von Interaktionsdatensätzen basiert, um die Interaktionsprognose zu optimieren.\vglfootcite[1]{Narayanan2023} Systeme wie FBLearner Flow von Facebook trainieren täglich tausende von Modellen an Billionen von Datenpunkten, um die persönliche Relevanz von Inhalten mit mathematischer Gewissheit vorherzusagen.\vglfootcite[18]{Zuboff2019} Social-Media-Plattformen verwandeln sich somit in gigantische Sensoren für gesellschaftliche Stimmungen und individuelles Verhalten, wobei die algorithmische Verarbeitung sicherstellt, dass die Informationsflüsse exakt auf die psychologischen Profile der Nutzer zugeschnitten werden.\vglfootcite[1]{VanDijck2014}\vglfootcite[2]{Lyon2012}

\subsection{Plattformökosysteme}

Die organisatorische und technologische Einbettung dieser Monitoring-Prozesse erfolgt innerhalb von Plattformökosystemen, die durch eine spezifische ökonomische Logik gekennzeichnet sind. Plattformen agieren als mehrseitige Märkte, die Interaktionen zwischen Nutzern, Werbetreibenden und Anbietern von Zusatzdiensten (Apps) koordinieren.\vglfootcite[664]{Cusumano2019}\vglfootcite[1]{Tiwana2014} Ein entscheidendes Element ist hierbei das Auftreten von Netzwerkeffekten: Je mehr Nutzer eine Plattform verwenden, desto wertvoller wird sie für alle Beteiligten, was zu einer natürlichen Tendenz zur Monopolisierung und einer immer umfassenderen Datensammlung führt.\vglfootcite[20]{Srnicek2017}

Innerhalb dieser Ökosysteme wird Daten zur strategischen Ressource, die den Kern der wirtschaftlichen Wertschöpfung bildet. Die Plattform-Governance legt dabei die Regeln fest, nach denen Daten extrahiert und genutzt werden dürfen, wobei die Plattformbetreiber als zentrale Torwächter fungieren.\vglfootcite[3]{Parker2016}\vglfootcite[1]{Tiwana2014} Durch die Integration verschiedener Dienste -- wie Suche, soziale Netzwerke und Cloud-Infrastrukturen -- können Unternehmen wie Amazon oder Google Daten aus unterschiedlichen Lebensbereichen konsolidieren, was die Vorhersagekraft ihrer Systeme massiv erhöht.\vglfootcite[51]{Srnicek2017}\vglfootcite[712]{Cusumano2019} Diese Ökosysteme schaffen einen technologischen Rahmen, in dem die Sammlung von Verhaltensdaten nicht mehr als Nebenprodukt, sondern als primärer Zweck der Infrastruktur erscheint, um durch algorithmische Optimierung garantierte kommerzielle Ergebnisse zu erzielen.\vglfootcite[4]{Zuboff2019}\vglfootcite[3]{Parker2016}

\subsection{Digitale Fußabdrücke und Nutzungsprofile}

Das Resultat der kontinuierlichen Datenerfassung ist die Transformation von digitalen Fußabdrücken (Digital Footprints) in detaillierte Nutzungsprofile. Jede Online-Aktivität hinterlässt \enquote{virtuelle Brotkrumen}, die ein präziseres Bild der Persönlichkeit und des Lebens eines Individuums zeichnen können, als dies durch bewusste Selbstauskünfte möglich wäre.\vglfootcite[28]{Zuboff2019} Monitoring-Systeme aggregieren diese Spuren zu umfassenden Datensätzen, wobei zwischen freiwillig bereitgestellten, beobachteten und erschlossenen Daten (Inferred Data) unterschieden werden muss.\vglfootcite[375]{OECD2011} Durch komplexe Daten-Mining-Verfahren und statistische Korrelationen können Systeme sensible Merkmale wie sexuelle Orientierung, politische Gesinnung oder die psychische Verfassung aus scheinbar banalen Interaktionen wie \enquote{Likes} ableiten.\vglfootcite[1]{VanDijck2014}\vglfootcite[1]{Acquisti2016}

Diese datengetriebene Repräsentation von Individuen führt zur Erstellung von digitalen Dossiers oder algorithmischen Identitäten, die oft als \enquote{Schattenkörper} (Shadow Bodies) bezeichnet werden, da sie nur jene Aspekte einer Person betonen, die für die algorithmische Logik lesbar sind.\vglfootcite[1]{Gillespie2014} Die Aggregation von Daten über verschiedene Plattformen hinweg ermöglicht es, Verhaltensmuster über lange Zeiträume zu modellieren und Individuen in immer feingranularere Segmente zu klassifizieren.\vglfootcite[1]{Kitchin2014}\vglfootcite[2]{Lyon2012} Diese Profile dienen nicht nur der Analyse der Vergangenheit, sondern bilden das Fundament für die Vorhersage und gezielte Beeinflussung zukünftigen Verhaltens durch algorithmische Steuerungsmechanismen.\vglfootcite[86]{Zuboff2019}\vglfootcite[379]{OECD2011}

Die systematische Erfassung von Verhaltensdaten und deren Verdichtung zu präzisen Nutzungsprofilen stellt jedoch nur die erste Phase des \ac{KI}-gestützten Monitorings dar. Sobald ein stabiles Abbild des Nutzers existiert, greifen die im nächsten Kapitel behandelten algorithmischen Steuerungs- und Bewertungsmechanismen, die durch Empfehlungssysteme, Rankings und Scoring-Verfahren aktiv in den Informationsfluss eingreifen, um die Interaktion zu maximieren und Handlungen gezielt zu lenken.

\section{Algorithmische Steuerungs- und Bewertungsmechanismen}

Die systematische Erfassung von Verhaltensdaten bildet die infrastrukturelle Grundlage für die operative Ebene \ac{KI}-gestützter Monitoring-Systeme, die durch spezifische algorithmische Verfahren eine aktive Steuerung von Informationsflüssen und Nutzerhandlungen vollziehen. Diese Mechanismen fungieren als instrumentelle Logik, die nicht nur Informationen filtert, sondern gezielt Relevanz konstruiert und so die digitale Realität der Nutzer strukturiert.\vglfootcite[1]{Gillespie2014}\vglfootcite[5]{Zuboff2019} Im Folgenden werden die technischen Funktionsweisen dieser Steuerungsmechanismen sowie deren Konsequenzen für die individuelle Verhaltenssteuerung erläutert.

\subsection{Personalisierung}

Die technische Realisierung einer individualisierten Informationsumgebung erfolgt primär durch Empfehlungssysteme (Recommender Systems), deren Kernaufgabe in der Schätzung des Nutzwerts (Utility) von bisher ungesehenen Objekten für einen spezifischen Nutzer besteht.\vglfootcite[1]{Adomavicius2005} Der Prozess basiert auf der Extrapolation bekannter Präferenzen in einen multidimensionalen Raum, wobei zwei grundlegende Paradigmen dominieren: Das inhaltsbasierte Filtern (Content-based Filtering) empfiehlt Objekte, die Ähnlichkeiten zu früher präferierten Inhalten aufweisen, während das kollaborative Filtern (Collaborative Filtering) Übereinstimmungen zwischen den Bewertungsmustern verschiedener Nutzer identifiziert, um Vorhersagen auf Basis von \enquote{People-to-People}-Korrelationen zu treffen.\vglfootcite[1]{Burke2002}\vglfootcite[1]{Adomavicius2005} Hybride Systeme kombinieren diese Ansätze häufig mit wissensbasierten Methoden, um die Vorhersagegenauigkeit weiter zu steigern und technologische Hürden wie das Kaltstart-Problem zu überwinden.\vglfootcite[1]{Burke2002}

Die Konsequenz dieser technologisch geschaffenen Umgebungen ist eine zunehmende Einengung des Informationshorizonts. Pariser (2011) beschreibt dieses Phänomen als Filterblase (Filter Bubble), in der die algorithmische Selektion dazu führt, dass Individuen primär mit dem vertrauten \enquote{Ghetto des Einen} konfrontiert werden.\vglfootcite[4]{Pariser2011} Da die Identität des Nutzers die Medieninhalte prägt und diese wiederum die Überzeugungen des Nutzers bestärken, entsteht eine Feedback-Schleife, die als \enquote{You Loop} bezeichnet wird.\vglfootcite[4]{Pariser2011} In der Logik des Überwachungskapitalismus dient diese Personalisierung nicht mehr nur der Nutzerzufriedenheit, sondern der Gewinnung eines Verhaltensüberschusses, indem das Individuum in eine algorithmisch lesbare und damit vorhersagbare Identität gedrängt wird.\vglfootcite[4]{Zuboff2019}\vglfootcite[1]{Gillespie2014}

\subsection{Ranking-Algorithmen}

Während die Personalisierung den individuellen Zuschnitt definiert, bestimmen Ranking-Algorithmen die hierarchische Ordnung und damit die Sichtbarkeit von Informationen innerhalb digitaler Plattformen.\vglfootcite[2]{Gillespie2014} Diese Algorithmen fungieren als neue Torwächter, die den Fluss von Wissen nach spezifischen, oft proprietären Relevanzkriterien steuern.\vglfootcite[1]{Gillespie2014}\vglfootcite[1]{Pasquale2015} Ein zentraler Mechanismus moderner Plattformen ist hierbei die Interaktionsprognose, bei der Inhalte danach gewichtet werden, mit welcher Wahrscheinlichkeit ein Nutzer eine Interaktion (Likes, Kommentare, Shares) vollziehen wird.\vglfootcite[1]{Narayanan2023} Ein prominentes Beispiel ist das Meaningful Social Interaction (MSI)-Modell von Facebook, das durch maschinelles Lernen Wahrscheinlichkeiten für verschiedene Interaktionstypen berechnet und den News Feed entsprechend priorisiert.\vglfootcite[1]{Narayanan2023}

Die daraus resultierende algorithmische Priorisierung transformiert soziale Interaktionen in eine ökonomische Logik, in der Aufmerksamkeit zur zentralen Währung wird. Da Ranking-Algorithmen bestimmen, was für die Öffentlichkeit wahrnehmbar ist, konstituieren sie sogenannte \enquote{berechnete Öffentlichkeiten} (Calculated Publics), die politische Diskurse und gesellschaftliche Wahrnehmungen massiv beeinflussen können.\vglfootcite[1]{Gillespie2014} Die Konsequenz dieser Machtasymmetrie besteht darin, dass Plattformbetreiber durch die Manipulation der Sichtbarkeit -- etwa durch die algorithmische Herabstufung (Demotion) von Inhalten -- gezielte Verhaltenssteuerung betreiben können, ohne dass die zugrunde liegenden Bewertungskriterien für den Nutzer transparent sind.\vglfootcite[1]{Narayanan2023}\vglfootcite[1]{Gillespie2014}

\subsection{Reputations- und Scoring-Systeme}

Eine spezifische Form der algorithmischen Bewertung ist die Umwandlung von Verhaltensdaten in automatisierte Scores, die zur Klassifikation und Hierarchisierung von Individuen dienen.\vglfootcite[1]{Pasquale2015} Durch die Aggregation von digitalen Fußabdrücken, wie etwa Klickströmen, Standortdaten oder sozialen Kontakten, werden Profile erstellt, die beispielsweise die Kreditwürdigkeit oder Zuverlässigkeit bewerten.\vglfootcite[2]{ONeil2016} Solche e-Scores fungieren als Instrumente zur Reduktion von Unsicherheit für Unternehmen, indem sie Individuen in mathematische Wahrscheinlichkeitskategorien einteilen.\vglfootcite[519\,f.]{OECD2011}\vglfootcite[2]{ONeil2016}

Diese Scoring-Verfahren haben tiefgreifende Konsequenzen für die soziale Teilhabe und Gerechtigkeit. O'Neil (2016) bezeichnet solche Modelle als \enquote{Weapons of Math Destruction} (WMD), da sie oft auf undurchsichtigen Kriterien basieren, Vorurteile kodifizieren und für die Betroffenen lebensverändernde Folgen haben können -- etwa den Ausschluss von Finanzprodukten oder Arbeitsplätzen --, ohne dass Einspruchsmöglichkeiten bestehen.\vglfootcite[2]{ONeil2016} Die algorithmische Bewertung tarnt dabei diskriminierende Praktiken, wie das \enquote{Redlining} auf Basis von Postleitzahlen, als objektive mathematische Evidenz und etabliert so eine Form der algorithmischen Governance, die das Leben des Einzelnen permanent überwacht und bewertet.\vglfootcite[1]{Pasquale2015}\vglfootcite[2]{ONeil2016}

\subsection{Interaktionsoptimierung und Verhaltensprognosen}

Die Integration von Personalisierung, Ranking und Scoring mündet in der systematischen Optimierung der Interaktion, um Verhaltensprognosen zu kommerzialisieren und Handlungen gezielt zu lenken. Plattformen nutzen hierbei Mechanismen des Bestärkenden Lernens (Reinforcement Learning), bei dem Agenten kontinuierlich aus Nutzerreaktionen lernen, um jene Aktionen (Policies) zu identifizieren, die den erwarteten Reward -- in Form von Aufmerksamkeit oder Klicks -- maximieren.\vglfootcite[1]{SuttonBarto2018}\vglfootcite[1]{Narayanan2023} Unterstützt wird dieser Prozess durch psychologische Trigger des Fogg Behavior Models, die Motivation und Fähigkeit des Nutzers adressieren, um durch gezielte Reize (Sparks/Signals) ein gewünschtes Zielverhalten zu automatisieren.\vglfootcite[1]{Fogg2009}

Praktische Anwendungen demonstrieren die Wirksamkeit dieser Verhaltensmodifikation. Die Facebook-Studie zur \enquote{Emotional Contagion} bewies, dass die algorithmische Manipulation von emotionalen Inhalten im News Feed die Stimmungslage von Millionen Nutzern ohne deren Wissen verändern kann.\vglfootcite[20]{Zuboff2019}\vglfootcite[1]{Kramer2014} Das Beispiel von Cambridge Analytica verdeutlicht zudem, wie psychometrische Profile (OCEAN-Analyse) für politisches Micro-Targeting genutzt wurden, um Wahlentscheidungen durch maßgeschneiderte Botschaften zu beeinflussen.\vglfootcite[1]{Isaak2018}\vglfootcite[7]{Zuboff2019} Schließlich führen diese Mechanismen zur Verstärkung von Echo-Kammern, in denen homophile Cluster und die algorithmische Belohnung von Emotionalität die soziale Polarisierung vorantreiben und die individuelle Autonomie zugunsten eines \enquote{Instrumentarianismus} untergraben.\vglfootcite[1]{Cinelli2021}\vglfootcite[1]{MetzlerGarcia2023}\vglfootcite[30]{Zuboff2019}

\section{Vergleich unterschiedlicher Systemlogiken}

Die Funktionsweise \ac{KI}-gestützter Monitoring-Systeme folgt je nach politischem und ökonomischem Rahmen unterschiedlichen Logiken, die jedoch auf einer gemeinsamen technologischen Basis der massenhaften Datenauswertung beruhen. Während westliche Systeme primär marktgesteuerten Akkumulationszielen folgen, zielen staatlich integrierte Systeme auf die Steuerung gesellschaftlicher Compliance ab. Der folgende Abschnitt analysiert diese Systemlogiken und stellt sie einander vergleichend gegenüber.

\subsection{Plattformzentrierte Systeme westlicher Plattformunternehmen}

In westlich geprägten Wirtschaftsräumen operieren Monitoring-Systeme innerhalb einer Logik, die als Überwachungskapitalismus beschrieben wird und in der menschliche Erfahrung als kostenloses Rohmaterial für die Extraktion von Verhaltensdaten dient.\vglfootcite[4]{Zuboff2019} Plattformunternehmen wie Google, Meta (Facebook), YouTube, TikTok und Amazon fungieren hierbei als zentrale Vermittler, deren digitale Infrastrukturen darauf ausgelegt sind, einen kontinuierlichen Fluss an Verhaltensüberschüssen zu generieren.\vglfootcite[19]{Srnicek2017}\vglfootcite[8]{Parker2016} Diese Systeme nutzen Algorithmen zur Personalisierung, um Informationsumgebungen so zu gestalten, dass sie die individuelle Relevanz maximieren und dadurch die Verweildauer sowie die Interaktionsrate erhöhen.\vglfootcite[4]{Pariser2011}\vglfootcite[1]{Adomavicius2005}

Die zentralen Ziele dieser plattformzentrierten Monitoring-Logik sind Interaktion, Wachstum und Monetarisierung. Die Interaktionsprognose, also die Vorhersage der Interaktionswahrscheinlichkeit, bildet den Kern der algorithmischen Steuerung, da ein engagierter Nutzer häufiger zur Plattform zurückkehrt und somit mehr Werbeeinnahmen generiert.\vglfootcite[1]{Narayanan2023} Netzwerkeffekte verstärken diesen Prozess, da ein steigendes Nutzeraufkommen den Wert der Plattform erhöht und noch umfassendere Datensammlungen ermöglicht, was wiederum die Vorhersagekraft der Algorithmen für Werbetreibende verbessert.\vglfootcite[20]{Srnicek2017}\vglfootcite[17]{Parker2016} Diese ökonomische Dynamik führt zur Entstehung von \enquote{berechneten Öffentlichkeiten}, in denen Ranking-Algorithmen die Sichtbarkeit von Inhalten nach kommerziellen Kriterien priorisieren, was oft zu einer Verengung des Informationshorizonts in Form von Filterblasen führt.\vglfootcite[1]{Gillespie2014} Die Konsequenz dieser Logik ist eine instrumentelle Macht, die das Verhalten der Nutzer nicht durch Zwang, sondern durch subtile algorithmische Anstöße und psychologische Trigger im Sinne der Plattformökonomie optimiert.\vglfootcite[4]{Zuboff2019}\vglfootcite[1]{Fogg2009}

Diese marktgetriebene Ausrichtung auf individuelle Präferenzen findet ihr Gegenstück in Systemen, die Monitoring-Technologien zur Sicherung staatlicher Autorität und gesellschaftlicher Ordnung einsetzen.

\subsection{Staatlich integrierte Monitoring-Systeme}

Staatlich integrierte Monitoring-Systeme, wie sie insbesondere in China entwickelt werden, nutzen dieselben technologischen Werkzeuge wie westliche Plattformen, jedoch unter einer fundamental anderen Zwecksetzung: der Sicherung von Compliance und staatlicher Governance. Das Herzstück dieser Entwicklung ist das Social Credit System (\ac{SCS}), das darauf abzielt, durch die Integration von Regierungsdaten und privaten Informationen ein umfassendes Maß an \enquote{Aufrichtigkeit} in wirtschaftlichen, sozialen und politischen Belangen zu erzwingen.\vglfootcite[1]{Creemers2018} Dieses System fungiert als kybernetischer Regelkreis, der Verhalten durch ein System von Belohnungen und Bestrafungen (Joint Punishment System) steuert, wobei Regelverstöße in einem Bereich zu Einschränkungen in allen Lebensbereichen führen können.\vglfootcite[1]{Creemers2018}\vglfootcite[1]{Qiang2019}

Technologisch stützt sich dieser digitale Autoritarismus auf eine omnipräsente Überwachungsinfrastruktur, die Künstliche Intelligenz mit massiven Videonetzwerken wie dem \enquote{Skynet}- oder \enquote{Sharp Eyes}-Projekt kombiniert.\vglfootcite[1]{Qiang2019}\vglfootcite[1]{SFRC2020} Regierungsbehörden setzen Verfahren wie Gesichtserkennung, Biometrie und Big-Data-Analysen ein, um Individuen in Echtzeit zu identifizieren, zu klassifizieren und Profile zu erstellen, die eine prädiktive Überwachung ermöglichen.\vglfootcite[1]{Qiang2019} Im Bereich des Smart Policing erlaubt die \enquote{Integrated Joint Operations Platform} (IJOP) die Aggregation verschiedenster Datenquellen, um verdächtige Verhaltensmuster automatisiert zu erkennen und präventive polizeiliche Maßnahmen einzuleiten.\vglfootcite[4]{SFRC2020}\vglfootcite[1]{Qiang2019} Diese Form der algorithmischen Governance zielt darauf ab, gesellschaftliche Risiken zu minimieren und die Stabilität des Regimes durch eine lückenlose Erfassung des Bürgerverhaltens sicherzustellen.\vglfootcite[1]{Creemers2018}\vglfootcite[2]{Lyon2012}

Der Vergleich dieser staatlichen Kontrolllogik mit der westlichen Plattformlogik verdeutlicht sowohl strukturelle Ähnlichkeiten als auch fundamentale Differenzen in der Zielsetzung und Machtausübung.

\subsection{Vergleich westlicher Plattformlogiken und staatlicher Monitoring-Systeme}

Ein Vergleich westlicher plattformzentrierter Systeme mit staatlich integrierten Monitoring-Systemen offenbart eine technologische Konvergenz bei gleichzeitiger Divergenz der zugrunde liegenden Governance-Strukturen. Beide Systemtypen basieren auf der massenhaften Sammlung digitaler Fußabdrücke und deren Transformation in automatisierte Scoring- und Klassifikationssysteme, sei es zur Ermittlung der Kreditwürdigkeit (e-Scores) oder der sozialen Vertrauenswürdigkeit (Social Credit).\vglfootcite[1]{Pasquale2015}\vglfootcite[1]{Creemers2018} Ebenso nutzen beide Seiten Mechanismen der Verhaltensmodifikation; während westliche Plattformen Nutzer durch Anstöße zu Konsumhandlungen bewegen, setzen staatliche Systeme auf eine offene Paternalisierung zur Einübung staatsbürgerlicher Tugenden.\vglfootcite[16]{Zuboff2019}\vglfootcite[1]{Creemers2018}

Die wesentlichen Unterschiede liegen in den Zielsetzungen und den Adressaten der Systeme. Westliche Logiken sind auf Interaktion vs. Compliance ausgerichtet: Das Ziel ist die Maximierung der Nutzeraktivität zur Monetarisierung von Aufmerksamkeit in einer Werbeökonomie, während staatliche Systeme auf die Einhaltung von Normen und Gesetzen zielen.\vglfootcite[1]{Narayanan2023}\vglfootcite[1]{Creemers2018} Dies spiegelt sich im Gegensatz von Wachstum vs. Governance wider; Plattformen streben nach globaler Expansion durch Netzwerkeffekte, wohingegen staatliches Monitoring der administrativen Effizienz und Risikokontrolle dient.\vglfootcite[20]{Srnicek2017}\vglfootcite[1]{SFRC2020} Der Adressat wird entsprechend unterschiedlich konstruiert: als Nutzer (Konsument/Datenquelle) im Westen und als Bürger im staatlich integrierten Modell.\vglfootcite[71]{Zuboff2019}\vglfootcite[1]{SFRC2020}

Ein weiterer markanter Unterschied besteht in der Form der Überwachung und Transparenz. Westliche Systeme operieren oft als Schwarze Kästen, deren proprietäre Algorithmen geheim gehalten werden, um Wettbewerbsvorteile zu sichern.\vglfootcite[1]{Pasquale2015} Im Gegensatz dazu strebt das chinesische Modell eine Form der radikalen Transparenz an, bei der Informationen über die Vertrauenswürdigkeit von Individuen und Unternehmen teilweise öffentlich bekannt gemacht werden, um den sozialen Druck zur Konformität zu erhöhen.\vglfootcite[1]{Creemers2018} Dennoch verschwimmen die Grenzen zusehends, da westliche Sicherheitsbehörden verstärkt auf die Datenbestände privater Plattformen zugreifen, während staatliche Systeme zunehmend Daten aus dem Privatsektor integrieren, was zu einem hybriden Überwachungskontext führt, in dem ökonomische und staatliche Steuerungslogiken ineinandergreifen.\vglfootcite[2]{Lyon2012}\vglfootcite[1]{Creemers2018}
